\begin{table*}[tbp]
\caption{RQ3 feasibility: anchor-associated endpoint changes under the example and word anchors.}
\label{tab:rq3-feasibility-anchor-patterns}
\begin{tabular}{lcc}
\toprule
Model & Example $\Delta_{\mathrm{endpoint}}$ [95\% CrI] & Word $\Delta_{\mathrm{endpoint}}$ [95\% CrI] \\
\midrule
Gemma3:12B & +0.001 [-0.005, 0.030]; Uncertain & +0.001 [-0.005, 0.031]; Uncertain \\
LLaMa-Pro & -0.083 [-0.161, 0.018]; Uncertain & +0.104 [-0.038, 0.275]; Uncertain \\
Mistral & +0.001 [-0.005, 0.035]; Uncertain & +0.006 [-0.004, 0.066]; Uncertain \\
Phi-4 & +0.001 [-0.005, 0.030]; Uncertain & +0.001 [-0.005, 0.031]; Uncertain \\
\bottomrule
\end{tabular}
\vspace{2pt}
\begin{minipage}{\linewidth}
\footnotesize\textit{Note.} $\Delta_{\mathrm{endpoint}}$ is the posterior anchor-minus-Context difference in $P(Y=1)+P(Y=4)$. Positive values indicate greater endpoint concentration under the anchor, whereas negative values indicate reduced endpoint concentration. A 95\% credible interval containing zero was classified as uncertain.
\end{minipage}
\end{table*}
